We study when explicit workflow information, defined here as a surgery's
overall phase order, improves remaining-surgery-duration (RSD)
prediction from surgical video. We compare a video-only baseline against
models that also use workflow information on two deliberately
contrasting benchmarks: MultiBypass140, a multicenter Roux-en-Y gastric
bypass dataset with meaningful workflow diversity, and Cholec80, a more
standardized single-center cholecystectomy dataset. Under a strict
prefix-only protocol that uses only frames available up to the
prediction time, workflow information improves MultiBypass140
within-center MAE on the development fold from 13.03 $\pm$ 0.18 to 12.18 $\pm$
0.11 min (3 seeds); the 5-fold $\times$ 3-seed mean gain is $-$0.19 min. In
cross-center evaluation, where the model is trained on Bern and
evaluated on Strasbourg, the same approach improves MAE from 17.80 $\pm$
0.55 to 17.33 $\pm$ 0.15 min. When we randomly reassign workflow labels
across surgeries, the gain disappears, showing that the improvement
comes from correct workflow information rather than from simply
providing extra conditioning information. A deployable predictor that
derives the workflow signal from the model's own phase predictions on
the observed prefix, rather than from the full case, still beats the
video-only baseline by 0.18 min within-center and 0.22 min cross-center
under a matched metric. On Cholec80, however, workflow information
provides little value: under the strict protocol, true workflow performs
similarly to the video-only baseline, and under the legacy
centered-window protocol prefix-based workflow estimates make prediction
worse. As a secondary absolute-performance result, our strongest
Cholec80 inference-time stack reaches 3.56 min MAE on the 30-video
public phase-labeled subset. This is lower than older published Cholec80
RSD results such as RSDNet (about 8 min) and TransLocal (7.10 min), but
it is not a like-for-like SOTA claim because it is measured on a
non-canonical subset, and the gain over our 4.46 min single-model
baseline comes mainly from isotonic calibration and 3-seed ensembling
rather than from workflow conditioning. Overall, workflow conditioning
helps when the benchmark contains real workflow variation, but provides
little benefit on more standardized benchmarks such as Cholec80.